\section{Method}\label{sec:method}

\subsection{Pre-registration}\label{sec:prereg}

We fixed the RQ1 design, the scorer and the RQ1 decision rule before held-out
set~1 was scored; the other designs and RQ2's confirmatory rule came later. One exposure came earlier: while designing M1 we viewed the
parsed tree structure of most documents in the first part of held-out set~1 (no
questions, answers or scores). Held-out set~1 was scored once, on 6~November,
and its results were reused (RQ5(b), answer quality, further reports and
checks); no chunker in the choice sets, scorer or decision rule was changed after
it was opened. A second held-out set, added later as a replication, was scored
once per gold variant, on 24~November 2026. The timeline, all deviations and every reuse are in
the repository (tag paper-v1.4).

\subsection{Documents and questions}\label{sec:docs}

\textbf{Tuning.} \TuningDocuments{} documents, of which \TuningMainDocuments{}
(\TuningMainQuestions{} questions) enter the main analysis, from
CUAD~\cite{hendrycks2021cuad}, OHR-Bench~\cite{zhang2024ocr}, an Indian
regulatory benchmark~\cite{singh2026parser}, QASPER~\cite{dasigi2021qasper} and
UniDoc-Bench~\cite{peng2025unidoc} (the last two left out).

\textbf{Held-out set 1.} \AccountingHeldOut{} sealed documents from known
sources, new document types and out of domain (one from
FinanceBench~\cite{islam2023financebench}). \AccountingMain{} documents with
\RQoneQuestions{} questions can be scored.

\textbf{Held-out set 2.} ViDoRe~V3~\cite{loison2026vidore}, a source new to the study,
selected from metadata alone (two earlier designs were replaced before any of their
documents was drawn or opened), keeping queries whose evidence pages are all text
or table pages: \HTwoQueries{} queries in \HTwoSelectedDocuments{} documents before
gold filtering. Under the primary gold rule, \RQoneHTwoDropRuleQuestions{}
questions in \RQoneHTwoDropRuleDocuments{} documents can be scored. 

\textbf{Question quality.} The OHR-Bench, UniDoc-Bench and Indian questions are LLM-generated. The author's check (one rater, stopped after its
first batch against its own expansion rule) found \GoldenOverallInvalid{} of
\GoldenOverallLabelled{} Indian tuning questions invalid.

\subsection{Gold}\label{sec:gold}

One matching rule places gold evidence as character spans in each parser's own
output. For held-out set~1 it is each source's answer or evidence span; a question that cannot be placed with certainty is dropped. This removes
\CoverageHeldOutDoclingShare{} of held-out set~1 questions for Docling and
\CoverageHeldOutPdfplumberShare{} for pdfplumber. For held-out set~2 the evidence
is the text inside the human-drawn boxes; the primary rule keeps a
page when the matcher's frozen drop rule accepts all its box text; two
sensitivity variants keep a query if any evidence page is placed  or
apply the held-out set~1 rule (certain-only). On 23~November, after seeing how many questions each rule keeps but before any
held-out set~2 score, we made the drop rule primary. Comparisons use only questions kept for every compared setup, which hides
coverage differences between parsers.

\subsection{Setups}\label{sec:setups}

A setup pairs a parser with a chunker setting. The parsers are
Docling~2.143.0~\cite{auer2024docling} (D, default settings) and
pdfplumber~0.11.10~\cite{pdfplumber} (P). The chunkers are fixed-token (FT),
recursive (Rec), LangChain recursive~\cite{langchain} (LC), sentence groups (Sent;
Stanza~\cite{stanza}), semantic (Qwen3-Embedding-0.6B~\cite{qwen3embedding};
tuning only, out of GPU memory on held-out data), LLM-regex
(Gemini~3.5 Flash-Lite~\cite{gemini}), page-based, and our structure-tree method, M1; four reuse the code of Adaptive
Chunking~\cite{demourajunior2026adaptive}. In names, the number is the size; +p marks a prefix,
LC-def LangChain's default and D-Tess Docling with Tesseract OCR. Chunks under 100 tokens are merged and those over 1,024 tokens split again. M1 keeps a parser-tree node whole if it fits the target
size, opens it otherwise, and packs neighbouring pieces.

The primary choice set was fixed before held-out set~1 was opened: Rec, LC, FT and M1 without prefix, each at 256, 384 and
512 tokens, plus sentence groups of five. That gives \RQonePrimarySettings{}
settings and, with two parsers, \RQoneSetups{} setups. Further choice sets are the original grid (\RQoneOriginalSetups{} setups on held-out set~1) and
the full grid (\HeldOutSetups{}).

\subsection{Retrieval and scoring}\label{sec:scoring}

Retrieval is BM25~\cite{bm25}, dense (bge-m3~\cite{chen2024m3}) or a hybrid by
reciprocal rank fusion~\cite{rrf}.
Recall@$B$ is the share of gold characters, each counted once, in the ranked
chunks read until $B$ tokens are used, the last one cut at the budget. We use $B \in \{750, 1{,}250, 2{,}000\}$; the headline result
is hybrid Recall@1,250. A floor test (chunks cut at random positions score \SanityFloorLower{} lower) failed
in all \SanityFloorFailedOriginal{} parser and retriever cases under its
pre-written rule and passed only under a rule amended after that result.

RQ3 (Recall@$k$, Hit@1 and MRR at fixed $k$ on tuning data) was designed on
12~November, after the held-out set~1 results, and has no test. A later control, also designed after
the held-out set~1 results, adds Lift@5.

\subsection{Statistics}\label{sec:stats}

Question scores are averaged per document, with 95\% intervals from a document
bootstrap~\cite{bootstrap} (2,000 document resamples, seed 42, two-sided). Two setups are tied when the paired interval of their per-document
difference contains zero. The tie rule changed twice before held-out set~1 was scored.

With split-half cross-fitting (the idea of~\cite{crossfitting}, not their
estimator) over 200 random splits, the oracle picks the best setup on half of a document's questions and
is measured on the other. The gap is the mean per-document recall of the oracle minus the best single setup. A document
with one question takes the best single setup's value. This holds for
\RQoneOneQuestionDocuments{} of \RQoneDocuments{} documents on held-out set~1 and
for \RQoneHTwoOneQuestionDocuments{} of \RQoneHTwoDropRuleDocuments{} on held-out
set~2, where the median document has \RQoneHTwoMedianQuestions{} questions. Choosing the best single setup on the questions it is scored on, and fitting the
oracle on half of a document's questions, both push the gap down, so it is not an
upper bound on what any chooser could gain. The rule: a per-document chooser is worth building only if the gap's interval lies
above zero. Multiple comparisons use Holm's correction~\cite{holm}.

RQ2 began as a descriptive tuning analysis, added on 12~November after the
tuning scores were known. A confirmatory rule was fixed before held-out set~2
was scored: RQ2 holds if, for each parser and budget, Spearman's
$\rho$~\cite{spearman} of measured chunk size against recall over that parser's
\RQtwoHTwoSetupsPerParser{} primary setups is negative. These setups span target
sizes 256 to 512 only.

\subsection{RQ4 and RQ5}\label{sec:aq}

RQ5(a), designed after held-out set~1 was scored, compares parser settings with defaults on digital tuning PDFs and scanned PDFs from
MP-DocVQA~\cite{tito2023mpdocvqa}. RQ5(b), whether a tuned chunk size beats the default, was specified on
10~November after the held-out set~1 scores existed: a second use of that set, not
confirmatory. 

For RQ4, the design, fixed after held-out set~1 was scored, had one
measure, an LLM judge; token F1 and
containment, added after the answers existed, are primary, and the judge (Qwen3-8B~\cite{qwen3}) failed its
human check and is secondary only. We sample \AQSampleShare{} of each PDF's questions and compare seven Docling
setups, and five scanned-PDF settings with Docling's default OCR or
Tesseract~\cite{tesseract}. Gemini Flash-Lite answers from the 1,250 tokens
hybrid retrieval reads.
